\reading{CR-19}{Future Value and Exposure}{DEFER}{3}

\begin{crkeybox}[Learning objectives for this reading]
\begin{enumerate}[label=\textbf{\alph*.},leftmargin=2em]
\item Describe and calculate the following metrics for credit exposure: expected mark-to-market, expected exposure, potential future exposure, expected positive exposure and negative exposure, effective expected positive exposure, and maximum exposure.
\item Compare the characterization of credit exposure to VaR methods and describe additional considerations used in the determination of credit exposure.
\item Identify factors that affect the calculation of the credit exposure profile and summarize the impact of collateral on exposure.
\item Identify typical credit exposure profiles for various derivative contracts and combination profiles.
\item Explain how payment frequencies and exercise dates affect the exposure profile of various securities.
\item Explain the general impact of aggregation on exposure, and the impact of aggregation on exposure when there is correlation between transaction values.
\item Describe the differences between funding exposure and credit exposure.
\item Explain the impact of collateralization on exposure and assess the risk associated with the remargining period, threshold, and minimum transfer amount.
\item Assess the impact of collateral on counterparty risk and funding, with and without segregation or rehypothecation.
\end{enumerate}
\end{crkeybox}

\begin{crtrapbox}[Single-layer reading, and the vocabulary reading for CVA]
CR-19 exists only in the Crash layer. It is also where the exposure metrics that
CR-20's CVA calculations consume are defined, so the definitions here are load
bearing further on.
\end{crtrapbox}

% =====================================================================
\lo{CR-19 a}{Describe and calculate the following metrics for credit exposure: expected mark-to-market, expected exposure, potential future exposure, expected positive exposure and negative exposure, effective expected positive exposure, and maximum exposure.}

\begin{center}
\begin{tikzpicture}
\begin{axis}[
  width=12.6cm, height=6.2cm,
  axis lines=left, xlabel={Future mark-to-market value},
  ylabel={Probability},
  xlabel style={font=\sffamily\scriptsize}, ylabel style={font=\sffamily\scriptsize},
  tick label style={font=\sffamily\scriptsize},
  xmin=-3.4, xmax=3.6, ymin=0, ymax=0.56, ytick=\empty,
  xtick={-3,-2,-1,0,1,2,3}, clip=true]
\addplot[domain=-3.4:3.6, samples=200, draw=none, name path=curve]
  {exp(-x^2/2)/sqrt(2*pi)};
\path[name path=base] (axis cs:-3.4,0) -- (axis cs:3.6,0);
\addplot[FRMGold!55] fill between[of=curve and base,
  soft clip={domain=-3.4:0}];
\addplot[FRMGreen!45] fill between[of=curve and base,
  soft clip={domain=0:3.6}];
\addplot[domain=-3.4:3.6, samples=200, smooth, draw=FRMNavy, line width=1.1pt]
  {exp(-x^2/2)/sqrt(2*pi)};
\draw[FRMNavy, line width=0.9pt] (axis cs:0,0) -- (axis cs:0,0.42);
\node[above, font=\sffamily\scriptsize\bfseries, text=FRMNavy, fill=white,
  inner sep=1.5pt] at (axis cs:0,0.42) {Expected MtM};
\draw[-{Stealth[length=4pt]}, FRMGreen, line width=0.9pt]
  (axis cs:0.8,0.03) -- (axis cs:0.8,0.30);
\node[above, font=\sffamily\scriptsize\bfseries, text=FRMGreen, fill=white,
  inner sep=1.5pt] at (axis cs:1.05,0.30) {EE};
\draw[-{Stealth[length=4pt]}, FRMGold, line width=0.9pt]
  (axis cs:-0.9,0.03) -- (axis cs:-0.9,0.30);
\node[above, font=\sffamily\scriptsize\bfseries, text=FRMGold, fill=white,
  inner sep=1.5pt] at (axis cs:-1.15,0.30) {NE};
\draw[-{Stealth[length=4pt]}, FRMRed, line width=0.9pt]
  (axis cs:2.6,0.01) -- (axis cs:2.6,0.22);
\node[above, font=\sffamily\scriptsize\bfseries, text=FRMRed, fill=white,
  inner sep=1.5pt] at (axis cs:2.6,0.22) {PFE};
\node[font=\sffamily\scriptsize, text=FRMGold, fill=white, inner sep=1.5pt,
  align=center] at (axis cs:-2.25,0.16) {average of NE\\is ENE};
\node[font=\sffamily\scriptsize, text=FRMGreen, fill=white, inner sep=1.5pt,
  align=center] at (axis cs:1.95,0.45) {average of EE\\is EPE};
\node[font=\sffamily\scriptsize\bfseries, text=FRMGold] at (axis cs:-2.6,0.02) {$(-)$};
\node[font=\sffamily\scriptsize\bfseries, text=FRMGreen] at (axis cs:2.9,0.02) {$(+)$};
\end{axis}
\end{tikzpicture}
\end{center}
\figcap{All seven metrics live on one distribution: the distribution of the
transaction's mark-to-market value at a future date. Everything to the right of
zero is exposure to \emph{you}; everything to the left is exposure to the
\emph{counterparty}. The right tail feeds CVA and the left tail feeds DVA, which
is why the two adjustments at CR-12 b have opposite signs.}

\begin{enumerate}
\item \term{Expected MtM.} The average future value of a transaction. It can
differ from the current MtM due to the passage of time and to cash flows.
\item \term{Expected exposure (EE).} The potential loss if positive MtM occurs and
the counterparty defaults. It is \emph{larger} than the expected MtM, because it
conditions on the positive side only.
\item \term{Potential future exposure (PFE).} The worst-case MtM scenario at a
specific future point.
\item \term{Expected positive exposure (EPE).} The average EE over time.
\item \term{Negative exposure (NE).} The counterparty's perspective: negative
future values. \term{Expected negative exposure (ENE)} is the opposite of EPE.
\item \term{Effective EE and effective EPE.} Capture \term{rollover risk} in
short-term transactions.
\item \term{Maximum exposure.} The highest PFE over the life of the transaction.
\end{enumerate}

\begin{crfmlbox}[Where these feed through to pricing]
\[ \mathrm{DVA} = -\mathrm{LGD}_I \sum_{i=1}^{m} \mathrm{ENE}(t_i)\times
   \mathrm{PD}_I(t_{i-1}, t_i) \]
\[ \mathrm{CVA} = -\mathrm{LGD}_C \sum_{i=1}^{m} \mathrm{EPE}(t_i)\times
   \mathrm{PD}_C(t_{i-1}, t_i) \]
\vk{Subscript $I$ denotes the institution making the calculation and $C$ the
counterparty. DVA is driven by the institution's \emph{own} default probability
applied to the negative exposure; CVA by the \emph{counterparty's} default
probability applied to the positive exposure.}
\end{crfmlbox}

\subsection*{Effective EE and effective EPE}

\begin{center}
\begin{tikzpicture}
\begin{axis}[
  width=12.2cm, height=5.6cm,
  axis lines=left, xlabel={Time}, ylabel={Expected exposure},
  xlabel style={font=\sffamily\scriptsize}, ylabel style={font=\sffamily\scriptsize},
  xtick=\empty, ytick=\empty,
  xmin=0, xmax=10, ymin=0, ymax=1.25, clip=true,
  legend style={font=\scriptsize\sffamily, draw=FRMRule, fill=white,
    at={(0.5,-0.22)}, anchor=north, legend columns=-1, column sep=10pt}]
\addplot[domain=0:10, samples=160, smooth, draw=FRMNavy, line width=1.1pt]
  {1.02*sin(deg(x*3.1416/10))^1.15};
\addlegendentry{EE}
\addplot[draw=FRMRed, line width=1.1pt, sharp plot]
  coordinates {(0,0) (1,0.42) (2,0.72) (3,0.91) (4,1.01) (5,1.02) (6,1.02)
  (7,1.02) (8,1.02) (9,1.02) (10,1.02)};
\addlegendentry{Effective EE (running maximum)}
\addplot[draw=FRMGreen, line width=1pt, dashed, sharp plot]
  coordinates {(0,0.86) (10,0.86)};
\addlegendentry{Effective EPE}
\addplot[draw=FRMGrey, line width=1pt, dotted, sharp plot]
  coordinates {(0,0.65) (10,0.65)};
\addlegendentry{EPE}
\end{axis}
\end{tikzpicture}
\end{center}
\figcap{Effective EE is the \emph{running maximum} of EE, so it can never fall.
That is what makes it capture rollover risk: a short-dated trade whose EE decays
to zero would otherwise look harmless, when in practice it is rolled and the
exposure returns. Because effective EE sits at or above EE everywhere, effective
EPE necessarily sits at or above EPE.}

% =====================================================================
\lo{CR-19 b}{Compare the characterization of credit exposure to VaR methods and describe additional considerations used in the determination of credit exposure.}

\begin{center}
\footnotesize
\begin{tabular}{L{2.9cm} L{5.5cm} L{6.1cm}}
\toprule
 & \thead{VaR} & \thead{Credit exposure} \\
\midrule
Purpose & Risk management only & \textbf{Both pricing and risk management} \\
Horizon & Short-term risks & Various \emph{longer} timeframes \\
Trends and volatility & Does not account for long-term market trends or
  volatility & Does \\
Cash flows & Overlooks future cash flows and contractual changes &
  Includes these to assess future risks \\
Mitigants & Does not include netting and collateral &
  Incorporates these, despite the complexity \\
\bottomrule
\end{tabular}
\end{center}
\figcap{Five differences, and the first is the one that drives the rest. Because
credit exposure is used for \emph{pricing} as well as risk management, it has to
carry everything that affects value over the life of the trade, not just the
shape of a short-horizon loss distribution.}

% =====================================================================
\lo{CR-19 c}{Identify factors that affect the calculation of the credit exposure profile and summarize the impact of collateral on exposure.}

Credit exposure gets more complex with:

\begin{enumerate}[label=\alph*)]
\item longer time horizons, which bring higher uncertainty;
\item uneven cash flows, which increase risk;
\item combined risk factors, for example cross-currency swaps;
\item options and exercise decisions; and
\item collateral, which reduces risk but brings complexities like delays and
value fluctuations.
\end{enumerate}

% =====================================================================
\lo{CR-19 d}{Identify typical credit exposure profiles for various derivative contracts and combination profiles.}

\begin{center}
\begin{tikzpicture}
\pgfplotsset{prof/.style={width=4.4cm, height=3.2cm, axis lines=left,
  xtick=\empty, ytick=\empty, xmin=0, xmax=1, ymin=0, ymax=1.15,
  title style={font=\sffamily\scriptsize\bfseries, text=FRMNavy},
  xlabel={\scriptsize time}, xlabel style={font=\sffamily\scriptsize,
  at={(0.5,-0.02)}}}}
\begin{axis}[prof, name=a, title={Bonds and loans}]
\addplot[draw=FRMNavy, line width=1pt, sharp plot]
  coordinates {(0,0.8) (1,0.8)};
\addplot[draw=FRMGreen, line width=0.8pt, dashed, smooth, domain=0:1, samples=30]
  {0.8+0.18*x};
\addplot[draw=FRMRed, line width=0.8pt, dotted, smooth, domain=0:1, samples=30]
  {0.8-0.28*x};
\end{axis}
\begin{axis}[prof, name=b, at={(a.right of south east)}, anchor=left of south west,
  title={Swaps}]
\addplot[draw=FRMNavy, line width=1pt, smooth, domain=0:1, samples=60]
  {2.6*x*(1-x)^0.85};
\end{axis}
\begin{axis}[prof, name=c, at={(b.right of south east)}, anchor=left of south west,
  title={FX products}]
\addplot[draw=FRMNavy, line width=1pt, smooth, domain=0:1, samples=60]
  {0.95*x^0.62};
\end{axis}
\begin{axis}[prof, name=d, at={(a.below south west)}, anchor=above north west,
  yshift=-16pt, title={Long options}]
\addplot[draw=FRMNavy, line width=1pt, smooth, domain=0:1, samples=60]
  {0.92*x^1.5};
\end{axis}
\begin{axis}[prof, at={(d.right of south east)}, anchor=left of south west,
  title={Credit default swaps}]
\addplot[draw=FRMNavy, line width=1pt, smooth, domain=0:0.72, samples=40]
  {0.55*x^0.8};
\addplot[draw=FRMRed, line width=1pt, sharp plot]
  coordinates {(0.72,0.44) (0.72,1.0) (1,1.0)};
\node[font=\sffamily\scriptsize, text=FRMRed, anchor=west]
  at (axis cs:0.24,0.93) {default};
\end{axis}
\end{tikzpicture}
\end{center}
\figcap{The five characteristic shapes. Bonds and loans sit at roughly the
notional, rising with interest rate risk (dashed) or falling with prepayments
(dotted). Swaps \emph{peak}, because future payment uncertainty rises while the
number of remaining payments falls. FX products increase monotonically, driven by
FX volatility and a large final notional exchange. Long options rise until
exercise. A CDS rises with widening spreads and then jumps to its maximum on
default, set by the recovery rate.}

\begin{itemize}
\item \term{Bonds and loans.} Exposure is roughly equal to the notional value,
with potential increases due to interest rate risk (bonds) or potential decreases
due to prepayments (loans).
\item \term{Swaps.} Typically show a \emph{peak} shape. This results from
balancing future payment uncertainties against the diminishing risk of payments
over time.
\item \term{Foreign exchange products.} Often have monotonically increasing
exposures, due to high volatility in FX rates, long maturities, and significant
final notional payments. The primary risk comes from fluctuations in FX rates
rather than interest rates.
\item \term{Long options.} Exposure increases over time until the option is
exercised, with the exact shape changing based on its position: in, out, or near
the money.
\item \term{Credit default swaps.} Exposure increases initially due to widening
credit spreads, then reaches a maximum upon default, based on the recovery rate.
\end{itemize}

\begin{crtrapbox}[Wrong-way risk in credit derivatives]
The effect of wrong-way risk leads to considerable counterparty risk for credit
derivatives: the event that makes the CDS valuable is correlated with the event
that makes the protection seller unable to pay.
\end{crtrapbox}

% =====================================================================
\lo{CR-19 e}{Explain how payment frequencies and exercise dates affect the exposure profile of various securities.}

\subsection*{Payment frequencies}

\begin{itemize}
\item \term{Unequal payments.} When payments \emph{received} are more frequent
than payments made, the exposure profile shows a \emph{reduction} in exposure.
More frequent cash inflows help mitigate the impact of future uncertainties.
\item \term{Equal payments.} When payments are made and received at the same
frequency, the exposure profile is typically smoother and less volatile.
\end{itemize}

\subsection*{Exercise dates}

\begin{center}
\footnotesize
\begin{tabular}{L{3.2cm} L{11.3cm}}
\toprule
\thead{Before the exercise date} & \term{Swaptions} generally have \emph{higher}
  exposure than forward swaps, due to the potential for the option not to be
  exercised if it is not profitable \\
\thead{After the exercise date} & \term{Forward swaps} usually have \emph{higher}
  exposure, because a forward swap has a fixed obligation to make payments while
  the swaption might not be exercised \\
\bottomrule
\end{tabular}
\end{center}
\figcap{The ranking flips at the exercise date. That flip is the whole content of
this objective, and it is trivially easy to state backwards.}

% =====================================================================
\lo{CR-19 f}{Explain the general impact of aggregation on exposure, and the impact of aggregation on exposure when there is correlation between transaction values.}

Netting agreements allow offsetting positions between parties to reduce overall
exposure. This is effective when trades have opposite MtM values.

\begin{itemize}
\item \term{Positive correlation} reduces netting benefits, as similarly directed
trades have little or no offsetting potential.
\item \term{Negative correlation} provides the most benefit, as perfectly
offsetting trades lead to a 100\% netting benefit, that is, zero overall risk.
\end{itemize}

\begin{crfmlbox}[Netting factor]
\[ \text{netting factor} \;=\; \frac{\sqrt{n + n(n-1)\bar{\rho}}}{n} \]
\vk{$n$ = number of exposures; $\bar{\rho}$ = average correlation. The factor is
the ratio of net exposure to the sum of individual exposures, so a
\textbf{lower netting factor is better}.}
\end{crfmlbox}

\begin{center}
\footnotesize
\begin{tabular}{L{3.0cm} L{4.6cm} L{6.9cm}}
\toprule
\thead{$\bar{\rho}$} & \thead{Netting benefit} & \thead{Netting factor} \\
\midrule
$\bar{\rho} = 1$ & No netting benefit & $100\%$ \\
$\bar{\rho} = 0$ & Good netting benefit & Depends on the number of exposures, but
  lower: $1/\sqrt{n}$ \\
$\bar{\rho} = -1$ & Maximum netting benefit & $0\%$ \\
\bottomrule
\end{tabular}
\end{center}

\begin{crtrapbox}[Correction: the $\bar{\rho}=-1$ row holds only for two exposures]
The claim that $\bar{\rho}=-1$ gives a netting factor of zero is true when
$n = 2$, and \emph{undefined} for $n > 2$. At $n = 3$ and $\bar{\rho} = -1$ the
expression under the root is $3 + 3(2)(-1) = -3$, and the square root of a
negative number is not available. The reason is that an average pairwise
correlation of $-1$ is not attainable across more than two exposures.

\smallskip
The correct general statement: the netting factor reaches zero at
\[ \bar{\rho} = -\frac{1}{n-1} \]
which is the \emph{lowest attainable} average correlation for $n$ exposures. For
$n=2$ that is $-1$, which recovers the notes' row; for $n=10$ it is only
$-0.111$.
\end{crtrapbox}

\begin{center}
\begin{tikzpicture}
\begin{axis}[
  width=12.2cm, height=5.6cm,
  xlabel={Number of exposures $n$}, ylabel={Netting factor},
  xlabel style={font=\sffamily\scriptsize}, ylabel style={font=\sffamily\scriptsize},
  tick label style={font=\sffamily\scriptsize},
  xmin=1, xmax=50, ymin=0, ymax=1.08,
  grid=major, grid style={FRMRule!50, line width=0.3pt},
  legend style={font=\scriptsize\sffamily, draw=FRMRule, fill=white,
    at={(0.5,-0.30)}, anchor=north, legend columns=-1, column sep=8pt}]
\addplot[domain=1:50, samples=80, smooth, draw=FRMRed, line width=1.1pt]
  {sqrt(x+x*(x-1)*1.0)/x};
\addlegendentry{$\bar{\rho}=1$}
\addplot[domain=1:50, samples=80, smooth, draw=FRMGold, line width=1.1pt]
  {sqrt(x+x*(x-1)*0.5)/x};
\addlegendentry{$\bar{\rho}=0.5$}
\addplot[domain=1:50, samples=80, smooth, draw=FRMNavy, line width=1.1pt]
  {sqrt(x+x*(x-1)*0.25)/x};
\addlegendentry{$\bar{\rho}=0.25$}
\addplot[domain=1:50, samples=80, smooth, draw=FRMGreen, line width=1.1pt]
  {sqrt(x)/x};
\addlegendentry{$\bar{\rho}=0$}
\end{axis}
\end{tikzpicture}
\end{center}
\figcap{Netting factor against portfolio size, for four average correlations. Two
readings. At $\bar{\rho}=1$ the line is flat at 100\%: adding trades buys nothing,
whatever the size. At $\bar{\rho}=0$ the factor falls as $1/\sqrt{n}$. In between,
each curve flattens quickly toward $\sqrt{\bar{\rho}}$ --- so beyond roughly
twenty exposures, extra trades add almost no further netting benefit and the
benefit is governed almost entirely by the correlation.}

% =====================================================================
\lo{CR-19 g}{Describe the differences between funding exposure and credit exposure.}

\begin{crkeybox}[The one-line mapping]
Counterparty defaults $\Rightarrow$ positive exposure $\Rightarrow$ \term{funding
costs}.

\smallskip
The bank itself defaults $\Rightarrow$ negative exposure $\Rightarrow$
\term{funding benefit}.
\end{crkeybox}

\begin{center}
\footnotesize
\begin{tabular}{L{2.7cm} L{5.6cm} L{6.2cm}}
\toprule
 & \thead{Funding exposure} & \thead{Credit exposure} \\
\midrule
Value definition & \textbf{Objective}: the cost of obtaining funds &
  \textbf{Subjective}: assumptions about loss recovery \\
Margin period of risk & Does not necessarily assume default &
  Considers default and the time to liquidate collateral \\
Aggregation & Considers the \emph{entire portfolio} for potential reuse of margin &
  Focuses on netting \emph{individual counterparties} \\
Wrong-way risk & Not considered &
  Considers factors that increase default likelihood while also decreasing
  collateral value \\
Segregation & Rules might be \emph{less} restrictive &
  Rules are more restrictive \\
\bottomrule
\end{tabular}
\end{center}

% =====================================================================
\lo{CR-19 h}{Explain the impact of collateralization on exposure and assess the risk associated with the remargining period, threshold, and minimum transfer amount.}

\begin{crdefbox}[Margin period of risk]
The MPoR is the time from a collateral call to the actual delivery of collateral.
It represents the critical exposure period for the party seeking collateral.
\end{crdefbox}

\begin{crfmlbox}[Exposure during the margin period of risk]
\[ \mathrm{EE} \;=\; 0.4 \times \sigma_E \times \sqrt{T_M}
   \qquad\qquad
   \mathrm{PFE} \;=\; k \times \sigma_E \times \sqrt{T_M} \]
\vk{$\sigma_E$ = annual volatility of the exposure; $T_M$ = the remargin period
expressed in years; $k$ = the quantile multiplier for the chosen confidence level.
For a 10-day remargin period, $\sqrt{T_M} = \sqrt{10/252} = 0.1992$.

\smallskip
The 0.4 is not arbitrary. For a zero-mean normal variable,
$E[\max(X,0)] = \sigma/\sqrt{2\pi} = 0.3989\,\sigma$, which rounds to 0.4. EE and
PFE are therefore the same quantity read at two different points of the same
distribution: its positive mean, and its $k$-th quantile.}
\end{crfmlbox}

\subsection*{Limitations of the PFE calculation}

\begin{enumerate}[label=\alph*)]
\item It assumes strong collateralization.
\item It does not account for uncertainty in collateral volatility. Collateral
volatility needs to be considered when using non-cash collateral.
\item It ignores liquidity and liquidation risks.
\end{enumerate}

\subsection*{Modelling collateral: which direction helps}

\begin{center}
\footnotesize
\begin{tabular}{L{5.4cm} L{9.1cm}}
\toprule
\thead{Parameter} & \thead{Direction that reduces exposure} \\
\midrule
Margin period of risk (MPoR) & \textbf{Lower} is better \\
Threshold & \textbf{Lower} is better \\
Minimum transfer amount & \textbf{Lower} is better \\
Initial margin & \textbf{Higher} is better \\
Rounding & --- \\
\bottomrule
\end{tabular}
\end{center}
\figcap{Four of the five have a direction, and initial margin is the odd one out:
it is the only parameter where \emph{more} is better. The other three are all
frictions that let exposure build before collateral arrives.}

\begin{crtrapbox}[Collateral is path-dependent]
The amount of collateral requested depends on how much was requested in the past.
The mechanism, through the minimum transfer amount, is set out at
\textbf{CR-17 c}.
\end{crtrapbox}

% =====================================================================
\lo{CR-19 i}{Assess the impact of collateral on counterparty risk and funding, with and without segregation or rehypothecation.}

\begin{center}
\footnotesize
\begin{tabular}{L{4.4cm} L{4.9cm} L{5.2cm}}
\toprule
\thead{Collateral form} & \thead{Counterparty risk} & \thead{Funding} \\
\midrule
\term{Cash, non-segregated} & Reduced & Reduced costs \\
 & \multicolumn{2}{l}{\emph{Best case: reduces both}} \\
\term{Securities, rehypothecated} & Mitigated if haircuts are sufficient &
  Reduced costs \\
 & \multicolumn{2}{l}{\emph{Good option: mitigates both}} \\
\term{Segregated cash or securities} & Reduced &
  \textbf{No funding benefit}, due to limited reusability \\
\term{Counterparty bonds, rehypothecated} &
  \textbf{Problematic}: they become worthless upon default &
  Reduced costs \\
\bottomrule
\end{tabular}
\end{center}
\figcap{The last two rows are the interesting ones and they fail in opposite
directions. Segregation protects against counterparty risk and kills the funding
benefit. Counterparty bonds deliver the funding benefit and are worthless exactly
when the counterparty risk crystallises, which is wrong-way risk inside the
collateral itself.}
